% TOP_PROBLEM: {"id": 20003265, "title": "Conditional fiber collapse and periodic-data tests for the Anosov infranil conjecture", "category_id": 20, "category": {"id": 20, "name": "miscellaneous", "display_name": "Miscellaneous", "description": "Problems whose source classification does not fit the main mathematical categories.", "slug": "miscellaneous", "order_index": 20, "created_at": "2026-07-31T15:26:25.671Z"}, "status": "partially_solved"}
% FIRST_POSTED: 2026-09-27
% ATTEMPT_STATUS: unresolved
% !TeX program = lualatex
\documentclass[11pt]{article}
\usepackage[margin=25mm]{geometry}
\usepackage{fontspec}
\setmainfont{Latin Modern Roman}
\usepackage{amsmath,amssymb,mathtools}
\usepackage{unicode-math}
\setmathfont{Latin Modern Math}
\usepackage{xurl,hyperref}
\hypersetup{colorlinks=true,urlcolor=blue}
\setcounter{secnumdepth}{0}
\setlength{\parindent}{0pt}
\setlength{\parskip}{5pt}
\setlength{\emergencystretch}{3em}
\begin{document}
\section*{\texorpdfstring{Conditional fiber collapse and periodic-data tests for the Anosov infranil conjecture}{Conditional fiber collapse and periodic-data tests for the Anosov infranil conjecture}}\label{20003265}
\subsection{Definitions and mathematical statement}
An Anosov diffeomorphism $f:M\to M$ of a closed smooth manifold has an invariant splitting $TM=E^s\oplus E^u$ and constants $C>0$, $0<\lambda<1$ with
$\|Df^nv^s\|\le C\lambda^n\|v^s\|$ and
$\|Df^{-n}v^u\|\le C\lambda^n\|v^u\|$ for $n\ge0$.
An infranilmanifold is a compact quotient of a simply connected nilpotent Lie group by a torsion-free discrete affine group with finite holonomy. The conjecture asks for a homeomorphism $h$ and a hyperbolic infranil automorphism $A$, with no eigenvalue of its inducing differential on the unit circle, such that
\[
 h\circ f=A\circ h.
\]
The precise infranil automorphism convention follows the source; topological conjugacy, not differentiable conjugacy, is requested.
\par\smallskip\textit{Formulation and concept references:} \href{https://aimath.org/WWN/measrigid/measrigid.pdf}{[S1]}.

\subsection{Short English statement}
Is every uniformly hyperbolic diffeomorphism topologically equivalent to an algebraic hyperbolic automorphism on an infranilmanifold?
\subsection{Research attempt}
\paragraph{Attribution and scope.}
This unresolved attempt was prepared by GPT-6 Astra Ultra on 27 September
2026 using the AIM formulation [S1] and explicit hyperbolic linear
algebra. The desired equivalence is topological conjugacy. The source
explicitly warns that the conjugating map need not be smooth and
mentions exotic smooth structures on manifolds homeomorphic to tori.
Those examples do not contradict the stated topological conclusion.
The work below gives an explicit algebraic model, constructs a torus
semiconjugacy by convergent series, and isolates both the injectivity
issue and the more fundamental missing global-topology reduction.
No classification of all Anosov manifolds is assumed.

\paragraph{A model with completely checkable dynamics.}
Let $A\in\mathrm{GL}(d,\mathbb Z)$ have no eigenvalue of modulus one.
It induces a diffeomorphism of $\mathbb T^d=\mathbb R^d/\mathbb Z^d$.
The real generalized eigenspaces split as $E^s\oplus E^u$, according
to eigenvalue modulus below or above one. Choose a number $\theta<1$
larger than all stable eigenvalue moduli and all unstable inverse
moduli. Jordan-block powers are a polynomial in the iterate times
an exponentially decaying factor. Absorbing that polynomial into a
slightly larger exponential rate yields constants $C$ such that
\[
 \|A^j|_{E^s}\|\le C\theta^j,
 \qquad \|A^{-j}|_{E^u}\|\le C\theta^j\quad(j\ge0).
\]
The constant distributions descend to invariant subbundles of the
tangent bundle, proving directly that the induced map is Anosov.
The unstable directions are not a global expansion in every direction;
the splitting is essential.

For a concrete surface model take
\[
                   A=\begin{pmatrix}2&1\\1&1\end{pmatrix}.
\]
Its eigenvalues are $(3\pm\sqrt5)/2$, with product one and one on
each side of one. This verifies all parts of the Anosov definition
without relying on a numerical Lyapunov exponent along selected orbits.

\paragraph{Periodic-point data give a necessary conjugacy test.}
For any positive integer $m$, a fixed point of $A^m$ corresponds to
$x\in\mathbb R^d$ with $(A^m-I)x\in\mathbb Z^d$, modulo
$\mathbb Z^d$. Since $A^m-I$ is invertible over $\mathbb R$,
\[
             \#\operatorname{Fix}(A^m)=|\det(A^m-I)|.
\]
This counts the finite cokernel of the integral matrix, equivalently
the index of its image lattice. Hyperbolicity guarantees invertibility
for every $m$, so no root-of-unity case is hidden in the formula.
For the displayed cat map, let $t_m=\operatorname{tr}(A^m)$.
Cayley--Hamilton gives
\[
 t_0=2,\quad t_1=3,\quad t_{m+1}=3t_m-t_{m-1},
 \qquad \#\operatorname{Fix}(A^m)=t_m-2.
\]
The first counts are $1,5,16,45,121$.

A topological conjugacy bijects fixed points at every iterate and
preserves least periods. Thus a proposed algebraic model must satisfy
these checks. Passing finitely many checks is insufficient, and even
all periodic counts would not construct a conjugacy. Conversely,
differences in derivatives at periodic points need not obstruct a
merely topological conjugacy; derivative multipliers are relevant to
stronger smooth-conjugacy questions. This keeps the diagnostic aligned
with the actual target.

\paragraph{An explicit semiconjugacy when the topology is already a torus.}
Assume now that $f:\mathbb T^d\to\mathbb T^d$ is a diffeomorphism
whose action on the fundamental group is a hyperbolic integral matrix
$A$. A lift has the form
\[
                   F(x)=Ax+\varphi(x),
\]
where $\varphi$ is continuous, bounded and $\mathbb Z^d$-periodic.
The lift is a homeomorphism. Let $P_s,P_u$ be the linear projections
onto the two invariant subspaces, and seek $H(x)=x+u(x)$ such that
$H\circ F=A\circ H$. This requires
\begin{equation}\label{a255:cohomology}
                    u(Fx)-Au(x)=-\varphi(x).
\end{equation}
The following uniformly convergent series solve it:
\[
 \begin{split}
 u_s(x)&=-\sum_{j=0}^{\infty}
                A_s^jP_s\varphi(F^{-j-1}x),\\
 u_u(x)&=\sum_{j=0}^{\infty}
                A_u^{-j-1}P_u\varphi(F^jx),
 \qquad u=u_s+u_u.
 \end{split}
\]
The exponential bounds and boundedness of $\varphi$ prove convergence
uniformly in $x$. Reindexing the stable series gives
$u_s(Fx)=A_su_s(x)-P_s\varphi(x)$, and reindexing the unstable series
gives the analogous identity. Their sum is
\eqref{a255:cohomology}. These sign checks matter: the stable part
uses backward iterates, whereas the unstable part uses forward iterates.

The function $u$ is periodic. Indeed
$F^j(x+z)=F^j(x)+A^jz$ for all integral $j$ and $z\in\mathbb Z^d$;
$A^jz$ remains integral, including negative $j$, because
$A\in\mathrm{GL}(d,\mathbb Z)$. Periodicity of $\varphi$ makes
every summand periodic. Therefore $H$ descends to a continuous map
$h:\mathbb T^d\to\mathbb T^d$ satisfying
\[
                             h\circ f=A\circ h.
\]
The maps $x\mapsto x+t u(x)$ descend to a homotopy from the identity
to $h$, so $h$ has degree one. It is surjective: if it omitted a point,
its degree would factor through the punctured torus, whose top
homology is zero. No claim that those homotopy maps are homeomorphisms
is used. The conclusion so far is semiconjugacy, not conjugacy.

\paragraph{A quantitative sufficient condition for injectivity.}
Suppose in addition that $f$ is expansive with constant $\delta>0$:
if $d(f^n x,f^n y)<\delta$ for every $n\in\mathbb Z$, then $x=y$.
Anosov diffeomorphisms have such constants by their local product
structure, an established dynamical fact. If
\[
                    2\sup_x d(h(x),x)<\delta,
\]
then $h$ is injective. If $h(x)=h(y)$, semiconjugacy gives
$h(f^n x)=h(f^n y)$ for every positive and negative iterate. The
triangle inequality bounds $d(f^n x,f^n y)$ by the displayed left
side, so expansivity forces equality. A continuous bijection of the
compact torus to itself is a homeomorphism, completing conjugacy under
this extra quantitative condition.

The series estimate supplies
\[
 \|u\|_\infty\le
 \left(\sum_{j\ge0}\|A_s^jP_s\|+
       \sum_{j\ge0}\|A_u^{-j-1}P_u\|\right)\|\varphi\|_\infty.
\]
Thus sufficiently small lift displacement relative to an available
expansivity constant gives a concrete criterion. This is not claimed
to cover all toral Anosov diffeomorphisms or to reprove the stronger
known rigidity theorems. For a large perturbation the bound may exceed
$\delta$, and bounded distance of lifted orbits is not the same as
the small distance required by expansivity on the torus.

\paragraph{Two obstructions to turning the calculation into a classification.}
First, a degree-one surjection need not be one-to-one. On the circle,
choose a continuous nondecreasing lift $g:\mathbb R\to\mathbb R$
with $g(t+1)=g(t)+1$, constant on a nontrivial interval and strictly
increasing elsewhere in a period. It descends to a degree-one
surjection that collapses an arc. Taking its product with identities
gives the same countertest on a torus. This example is not asserted
to be the semiconjugacy of an Anosov map; it shows that the degree
argument alone has not proved injectivity. Hyperbolicity must enter
again through an additional argument, such as the quantitative one
above or a stronger global rigidity theorem.

Second, the entire series construction assumed the manifold was a
torus and the induced matrix hyperbolic. A general Anosov splitting
does not supply a global integral affine coordinate system. Even
proving a particular manifold has the homotopy type of an infranilmanifold
would still require the appropriate topological identification and
dynamical conjugacy. The existence of contracting and expanding
subbundles is differential information; it does not directly prove that
the fundamental group is virtually nilpotent.

The statement also concerns diffeomorphisms, not time-one maps of
Anosov flows. Such a time-one map has a central flow direction and
fails the displayed two-subbundle hyperbolicity definition. Importing
nontransitive flow examples without checking that direction would
produce spurious counterexamples to the catalogue assertion.

\paragraph{Verification and outcome.}
Exact matrix powers check the periodic-point formula and recurrence.
The semiconjugacy is verified by reindexing absolutely uniformly
convergent series, including its lattice equivariance. The injectivity
step exposes its required smallness hypothesis rather than replacing
boundedness by expansivity. The strongest result here is an explicit
semiconjugacy construction and a sufficient conjugacy criterion within
an already algebraically modeled topology. The first general gap is
obtaining the infranil topological model from an arbitrary Anosov
system; no construction or counterexample resolves it. The full
classification problem remains \textbf{UNRESOLVED}.

\subsection{Sources}
\begin{itemize}
\item[S1]Open Problems from the Workshop 'Emerging Applications of Measure Rigidity' (AIM, 2004), Section 2.1, Conjecture 3.
\url{https://aimath.org/WWN/measrigid/measrigid.pdf}.
\textit{Formulation source.}
\end{itemize}
\end{document}
